% The appendices of the notes.  Chapter A is the method, written for the
% students; then one chapter per research question (B: RQ1, C: RQ2 and
% RQ3, D: RQ4).  Each research chapter ends with its author's block and
% the hit-list table.
%
% Author's note: drafted 2026-11-05 from three review reports by the
% agents lit-rq1, lit-rq23 and lit-rq4 (scholar sessions lp-llm-direct,
% lp-prose-effects, lp-agent-context).  Every key cited here is in
% ltnotes.bib with a provenance block; every \textcquote is copied from a
% QUOTE line of that block.

\chapter{Checking a claim}
\label{app:ai-method}

Several things this lecture says are claims about what research has found,
and you can't check them at the keyboard the way you can check what a
program prints.
This chapter says how we checked them, so that you can judge how far to
trust them and redo the checks yourself.
The chapters after it take the claims one research question at a time.

The lecture rests on the literature at these points:
\begin{description}
  \item[What has been tested] Since 2023, literate programming has been
    connected with language models mostly in proposals and tools, only four
    works measure anything, and none tests literate programming as Knuth
    meant it.
    This is the answer to \cref{ex:ai-research-amount}, with the works sorted
    by \cref{def:ai-evidence}.
    \Cref{app:ai-rq1} asks how the two have been connected since 2023, and
    finds that the claim holds, in a literature that is small and consists
    mostly of preprints.
  \item[Prose for the agent] In \cref{sec:ai-agent-reader}: prose helps a
    language model when it tells the model something correct that it
    couldn't otherwise know, wrong prose hurts, and prose that restates the
    code does little.
    \Cref{app:ai-rq23} finds this supported by benchmarks, with moderate
    confidence, though none of them uses a literate program.
  \item[Prose for the checker] In \cref{sec:ai-checker}: explanations next
    to the code helped people in some studies, but not in the one experiment
    that varied how good the explanations were; comments make people adopt
    AI-written code more often; and AI assistance can lower understanding.
    \Cref{app:ai-rq23} finds the evidence mixed, and weaker than for the
    model.
  \item[Prose that instructs the agent] In \cref{sec:ai-instructions} and
    with \cref{ex:ai-rate-instructions}: context files did not significantly
    change whether agents solved their tasks, at a higher cost; out-of-date
    files are followed over the code; skills written by people helped on
    specialized tasks but mostly not on ordinary repositories; and
    spec-driven development has hardly been evaluated beyond its authors'
    own projects.
    \Cref{app:ai-rq4} supports each of these, with the reservation that
    nearly all of the evidence is preprints from 2026.
  \item[Proof of agent work] With \cref{ex:ai-rate-showboat}: we found no
    scientific evaluation of documents like Showboat's.
    \Cref{app:ai-rq4} reports the closest studies.
\end{description}
The pattern the lecture draws across these studies, that prose helps an
agent when it carries information the agent can't get from the code or
doesn't already know, is our synthesis of \cref{app:ai-rq23,app:ai-rq4}.
No study tested it, and the lecture says so.

\section{What kind of review this is}

The searches follow the method of a systematic review: written research
questions, the same queries in every database, searches both for and against
each claim, and a record of every step.
They are not a full systematic review, though, and you should read them as
systematic in style only:
\begin{itemize}
  \item The screening was done in a single pass, by AI agents under the
    teacher's supervision, with a language model classifying the records.
  \item No second, independent reviewer checked the screening or the
    classification.
  \item All of it was done in one evening, on 5 November 2026.
  \item Most of the relevant work is from 2026, much of it preprints that
    have not been peer-reviewed, and it measures model versions that will
    soon be replaced.
\end{itemize}

\section{Turn the claim into a question}

A claim can be checked only once it's clear what would show it wrong.
\enquote{Documentation helps with agents} can't be shown wrong, because it
doesn't say whom it helps, to do what.
That is why the lecture separates the readers: the agent that works from the
prose, and the person who checks what the agent wrote.
Each becomes its own question, and each question says what counts as help:
whether the model solves the task, what it costs, or whether the person
judges the code correctly.

The question also fixes how strong the claim is.
\enquote{Can help} needs one good study; \enquote{helps} needs several that
agree, and no good study against it.
A claim that something has \emph{not} been done, such as \enquote{nobody has
tested this}, can be supported by a search but never proved by one, so the
chapters say \enquote{we found no} rather than \enquote{there is no}.

\section{Choose the right kind of source}

Which source can answer a question depends on the question.
What literate programming is, we take from the primary source, Knuth's
article\autocite{Knuth1984}, not from someone's summary of it.
Whether something works, a single study answers for its own setting, and a
review of many studies answers for whether the result is established.
For the questions of this lecture there are hardly any reviews yet, so the
chapters weigh single studies against each other.

Two more things decide how much weight a source gets.
The first is the kind of evidence, sorted as in \cref{def:ai-evidence}: a
proposal shows that something can be built, a benchmark shows what a model
does, and only a study with people shows what people do.
The second is the source's status: whether it has been peer-reviewed or is a
preprint that its authors have posted themselves, and whether the authors
sell the product they evaluate.
The chapters give both for every source they lean on.

\section{Search so that you can be wrong}

A search that looks only for support finds it.
Each search therefore has two kinds of query.
Support queries look for work on the claim.
Counter-queries use the same terms together with words such as
\enquote{limitations}, \enquote{does not improve} or \enquote{misleading}, to
find work that contradicts or narrows the claim.
The queries name concepts, not authors, since a search for the people you
expect finds the work you already believed in.

Every query was sent to the same databases: OpenAlex, Web of Science, IEEE
Xplore, arXiv and HAL, through the command-line search tool
\texttt{scholar}, which records each query and its hits.
Two further databases, Semantic Scholar and DBLP, failed on the day, so the
chapters mark them as not searched, never as having no hits.
Databases don't find everything, so we also followed citations, both to the
papers that cite Knuth's article and from the reference lists of the
central papers.
Searches outside the databases, in the paper search service alphaXiv and on
the web, are given in each chapter with their exact wording.

Before searching, we set aside a few papers known to be relevant, to test
the queries.
Where the queries missed one, we improved them with the terms the field
uses, rather than adding the missing paper by name, which would have hidden
how much else the queries were missing.

\section{Read the full text, and write down how}

A hit is not evidence.
A title can promise more than the paper shows, and an abstract can report
the best number without the caveat.
Each cited source was opened, and the statement it backs was matched against
a verbatim passage, which the chapters quote with its location.
Where only the abstract could be read, the chapter says so.

For each source, the source file of the bibliography records a fixed set of
fields:
\begin{description}
  \item[Claim] The statement the source is cited for.
  \item[Found via] The exact query, database or citation trail that found
    it.
  \item[Picked] Why this source rather than another.
  \item[Quote] The passage that supports the claim, verbatim, with its
    location.
  \item[Verified] What was read, full text or abstract only, and anything
    that limits how the source applies.
  \item[Counter] The counter-search and what it found.
  \item[Date] When the source was checked.
\end{description}

Every record the database searches returned was then classified, first by a
language model and then again by the reviewing agent for every record that
bears on the question, as supporting the claim, qualifying or contradicting
it, adjacent (the same field, another question) or a false hit.
The table at the end of each chapter lists the records that bear on the
question, so that you can see what was left out as well as what was cited.
For \cref{app:ai-rq23} the classification was not finished in time, so its
table is still to come.

\section{Draw the conclusion, with reservations}

The conclusion of each chapter starts by answering its question, at the
strength the evidence supports, and then says what limits the answer:
databases that failed, papers read only in abstract, results that rest on a
single study.
The lecture states each claim no more strongly than its chapter concludes.
Each chapter ends with further work: what its search left undone, and what
the field would have to study for a surer answer.

\section{A note on language models}

The searches, the screening, the classification and the first drafts of
these appendices were done by AI agents (Claude, working in Claude Code)
under the teacher's supervision, and the language model \texttt{gpt-5.6-sol}
classified the records.
That is itself an instance of what the lecture is about: prose about the
research, written by an agent, which a person has to check.
An agent can misread a paper or quote from the wrong version of it, so the
chapters give the passages verbatim and say where they are, which lets you
check them.
The agents were told to cite only what they had opened and to record what
they could not read.
As the note at the start of each chapter says, the teacher has not yet
reviewed the results in full.


% ================================================================ RQ1
\chapter{How has literate programming been connected with language models
  since 2023?}
\label{app:ai-rq1}
\chapterprecis{The author has not yet reviewed the results in this
  appendix in full.}

The lecture claims that since 2023 literate programming has been connected
with language models mostly in proposals and tools, that only four works
measure anything, and that nobody has tested literate programming as Knuth
meant it together with a language model (\cref{ex:ai-research-amount}, with
the evidence sorted by \cref{def:ai-evidence}).
By literate programming we mean what \textcite{Knuth1984} proposed:
\textcquote[p.~97]{Knuth1984}{explaining to human beings what we want a
computer to do}, in a program whose concepts are introduced
\textcquote[p.~97]{Knuth1984}{in an order that is best for human
understanding}.

The research question is: how has literate programming been connected with
large language models, generative AI or coding agents since 2023, and how
much of that work has been evaluated?
A work is in scope if it explicitly connects the two.
The claim has two parts of different strength.
That the work is mostly proposals is a statement about proportions, which
the search can test.
That no controlled test exists is a statement about absence, which the
search can support but not prove.

\section{Method}
\label{sec:ai-rq1-method}

% Author's note (verification and reproduction): see the block above the
% hit-list \input at the end of the chapter.

The search was made on 5 November 2026 with \texttt{scholar}.
Its check of the databases that evening found OpenAlex, Web of Science, IEEE
Xplore, arXiv and HAL working.
Semantic Scholar rejected the access key (HTTP 403), and DBLP answered with a
page that stops automated clients instead of with data (HTTP 202), so
neither could be searched.
Scopus answered, but was not part of the search plan and was not used.
Every query was sent to each working database separately, limited to works
from 2023 onwards and to 100 hits per database.
Web of Science received the same keywords as a topic search,
TS=(\dots).

The queries combine four groups of terms:
\begin{description}
  \item[LLM] LLM OR LLMs OR \enquote{large language model} OR
    \enquote{large language models} OR \enquote{generative AI} OR GenAI OR
    ChatGPT OR Copilot OR \enquote{code generation} OR \enquote{coding
    agent} OR \enquote{coding agents} OR \enquote{AI-assisted programming}
    OR \enquote{vibe coding}
  \item[LP2] \enquote{literate program} OR \enquote{literate programs} OR
    noweb OR CWEB OR \enquote{org-babel} OR \enquote{org-mode} OR (Knuth AND
    literate)
  \item[NB] \enquote{computational notebook} OR \enquote{computational
    notebooks} OR Jupyter
  \item[CTR] limitation OR limitations OR criticism OR critique OR failure
    OR \enquote{negative result} OR \enquote{negative results} OR
    \enquote{does not} OR overhead OR cost
\end{description}
The queries themselves are:
\begin{description}
  \item[S1] \enquote{literate programming} AND (LLM)
  \item[S2] (LP2) AND (LLM)
  \item[S3] (NB) AND literate AND (LLM)
  \item[S4] literate AND (LLM)
  \item[C1] \enquote{literate programming} AND (LLM) AND (CTR)
  \item[C2] (LP2) AND (LLM) AND (CTR)
  \item[C3] (NB) AND literate AND (LLM) AND (CTR)
  \item[C4] \enquote{literate programming} AND (adoption OR \enquote{lack
    of adoption} OR barrier OR barriers OR criticism OR critique OR
    limitations OR overhead OR cost)
  \item[U1] \enquote{literate programming}
\end{description}
The support queries S1--S4 look for the connection, from the exact phrase
(S1) to any use of \enquote{literate} (S4), which is broad on purpose and
catches the other meanings of the word.
The counter-queries C1--C3 add the CTR terms to S1--S3, to find
limitations and negative results, and C4 asks about the adoption problems of
literate programming itself.
U1 takes everything on literate programming since 2023.
\Cref{tab:ai-rq1-queries} gives the hits.

\begin{table}[htbp]
  \begin{sidecaption}{Queries and number of hits per database, 5 November
    2026, with the tool \texttt{scholar}; the queries are given in
    \cref{sec:ai-rq1-method}.
    S = support query, C = counter-query, U = catch-all query.
    Databases: OpenAlex (OA; not open access), Web of Science (WoS), IEEE
    Xplore (IEEE), arXiv, HAL (Hyper Articles en Ligne), Semantic Scholar
    (S2), DBLP.
    At most 100 hits were fetched per database and query, so 100 means
    \enquote{at least 100}.
    \enquote{--} means that the search failed: Semantic Scholar rejected the
    access key and DBLP answered every query with a page that stops
    automated clients, so neither database was searched.
    Web of Science got the keywords as \texttt{TS=(\dots)}, the others as
    boolean free text.}[tab:ai-rq1-queries]
    \centering\footnotesize
    \begin{tabular}{@{}lrrrrrrr@{}}
      \toprule
         & OA  & WoS & IEEE & arXiv & HAL & S2 & DBLP \\
      \midrule
      S1 & 100 & 3   & 2    & 4     & 0   & -- & --   \\
      S2 & 100 & 1   & 1    & 4     & 1   & -- & --   \\
      S3 & 100 & 0   & 0    & 0     & 0   & -- & --   \\
      S4 & 100 & 60  & 19   & 100   & 4   & -- & --   \\
      \midrule
      C1 & 100 & 0   & 0    & 1     & 0   & -- & --   \\
      C2 & 100 & 0   & 0    & 1     & 0   & -- & --   \\
      C3 & 100 & 0   & 0    & 0     & 0   & -- & --   \\
      C4 & 100 & 5   & 2    & 7     & 1   & -- & --   \\
      \midrule
      U1 & 99  & 31  & 7    & 13    & 12  & -- & --   \\
      \bottomrule
    \end{tabular}
  \end{sidecaption}
\end{table}

OpenAlex matches the full text, so most of its hits mention a term only in
passing, and every OpenAlex cell but one reached the limit of 100.
Its totals behind the limit, for the same queries, were 291 (S1), 371 (S2),
235 (S3), 11,505 (S4), 278 (C1), 358 (C2), 228 (C3), 567 (C4) and 759 (U1).
Because of that limit, the database searches were complemented the same day
by searches that do not depend on relevance ranking:
\begin{description}
  \item[Forward citations] All 180 works from 2023 onwards that cite
    Knuth's article, according to OpenAlex, were screened by title and
    abstract for terms about AI.
  \item[Title and abstract] An OpenAlex search for \enquote{literate
    programming} in the title or abstract, from 2023 onwards and without a
    limit, gave 92 works.
  \item[Backward citations] The reference lists of eight of the central
    papers were screened for literate programming work from 2023 onwards.
    They returned only works already found.
  \item[alphaXiv] Two searches, both for papers from 2023 onwards: one with
    the keywords \enquote{literate programming} and \enquote{LLM} and the
    question \enquote{literate programming with LLMs, generative AI or
    coding agents}, and one with the keywords \enquote{literate
    programming}, \enquote{computational notebook} and
    \enquote{limitations} and the question \enquote{limitations, criticism
    or negative results of literate programming or literate notebooks with
    LLM coding assistants}.
  \item[The web] Two searches: \enquote{literate programming coding agents
    LLM evaluation empirical study noweb org-mode 2025 2026}, and
    \enquote{literate programming AI agents criticism does not work overhead
    why literate programming failed adoption LLM}.
\end{description}
The alphaXiv and web searches added one source each, both cited in
\cref{sec:ai-rq1-results}.
Knuth's article was retrieved as a known document and read in the published
version, from the teacher's own copy, since the publisher refused the
download.

Before the search, eight papers known to be relevant were set aside to test
the queries.
The queries found seven of them.
The eighth, on \enquote{literate tracing}\autocite{Sotoudeh2025}, was
found only through alphaXiv; it never mentions language models, so it is
outside the question, and the queries were right to miss it.

The 1,178 database hits came to 567 unique records, and with the forward
citations and the title-and-abstract search, 729 unique works were screened.
A language model classified the records of the database searches as
connecting literate programming with language models explicitly, qualifying
that connection, adjacent, or false hits.
Every record it put in the first two groups, and every record it was unsure
of (a confidence under 0.7), was read again, and four of its decisions were
changed.
The other 149 works were screened by searching their titles and abstracts
for terms on both topics and reading the abstracts that matched.
In addition, 39 arXiv full texts among the hits of S1 and S2 were searched
for \enquote{literate programming} and Knuth: four use literate programming
centrally, and 22 mention it only in passing.
Eleven works are in scope: three peer-reviewed papers, one paper in
conference proceedings, six preprints and one blog post.

\section{Results}
\label{sec:ai-rq1-results}

\subsection{What has been proposed}

Most of the work proposes a way to use literate programming with language
models, without measuring whether it works.
At Google, \textcite{Shi2025NLO} let a model write \enquote{natural-language
outlines}, prose that \textcquote[abstract]{Shi2025NLO}{partition the code
and summarize its main ideas in the style of literate programming}.
\Textcite{Zhang2024ILP} use \enquote{interoperable literate programming}
documents as the prompt for a model working on a large project; their
evidence is a demonstration in which GPT-4o wrote correct code for all of
twelve small Scheme functions, and the authors develop the tools the paper
promotes.
\Textcite{Yuan2026VLP} propose \enquote{verifiable literate programming},
documentation derived from the code that a person checks instead of the
code itself.
\Textcite{Burke2026notlob} build a literate programming environment for
humans and agents, and state that \textcquote[sec.~1]{Burke2026notlob}{%
experimental evaluation is left to future work}.
\Textcite{VanClief2026ICM} describe a way of organizing an agent's context
in folders of Markdown files and note that \textcquote[sec.~3]{%
VanClief2026ICM}{There is also something of Knuth's literate programming in
this design}; their support is informal reports from the method's own
community.

\subsection{What has been measured}

Four works measure something, each in its own way.
The only journal article, by \textcite{SunStaron2026}, asks whether the
internal representations of 19 models recognize both the programming
language and the task of a program, and concludes that
\textcquote[sec.~VIII]{SunStaron2026}{trillion-scale models are the only
ones that come close to literate programming by consistently recognizing
both the programming language and the programs' corresponding tasks with an
accuracy of over 70\,\% on both datasets}.
The sizes of those models are the authors' estimates, and the authors
themselves note that the study \textcquote[sec.~VII]{SunStaron2026}{%
evaluates \enquote{literacy} indirectly through classification accuracy
rather than through end-to-end code generation or explanation}.
No literate program is written or read in it.

In the outline study, six developers rated the outlines of 30 functions they
had chosen from their own projects; five of the six raters are the paper's
authors, and \textcquote[sec.~5]{Shi2025NLO}{60\,\% of
outlines were rated as having excellent overall quality (90\,\% were
acceptable or better)}; in a case study, 30 Android reverse engineers rated
83.8\,\% of the outlines completely correct.
The authors note the small samples themselves, and the study rates the
prose, not whether it helps anyone.

The verifiable literate programming preprint reports that it
\textcquote[abstract]{Yuan2026VLP}{improves code pass@1 from 28.7\,\%-73.2\,\%
to 65.4\,\%-93.5\,\%}, where pass@1 is the share of tasks solved at the first
attempt.
Those numbers come from a language model playing the user, and
\textcquote[sec.~V-A]{Yuan2026VLP}{The oracle context includes the
benchmark prompt and golden reference implementation}, so the simulated
user knows the answer.
The study with people had eight participants.

The one study that isolates literate comments is a benchmark of
mathematical modelling by \textcite{RossiPrestwich2026}, in which the
\textcquote[sec.~5]{RossiPrestwich2026}{ablation of literate-style comments
\textelp{} did not produce differences in accuracy with respect to the
baseline}.
The authors' further claim, that the comments make the models more readable,
rests on their own inspection, not on a measurement.

\subsection{Positions and practice}

Two education papers take positions, read only in abstract since the full
texts could not be obtained.
\Textcite{Birkenkrahe2024} posits that AI coding assistants and literate
programming are complementary, a \textcquote[abstract]{Birkenkrahe2024}{%
symbiotic relationship}.
\Textcite{Rosiene2024} argue the opposite for teaching: \textcquote[%
abstract]{Rosiene2024}{In contrast to the practice of Literate Programming},
students should learn to read LLM-generated code in a way that
\textcquote[abstract]{Rosiene2024}{requires little or no commentary}.
A practitioner's essay, found by the web search, argues that
\textcquote{Whitlock2026}{The fundamental extra labor of literate
programming, which I believe is why it is not widely practiced, is
eliminated by the agent}, and adds \textcquote{Whitlock2026}{I don't have
data to support this}.
None of these has data, so their disagreement is one of positions, not of
findings.

\subsection{Neighbouring work}

Two works continue the literate programming tradition without language
models.
\Textcite{Bond2026LitExec} design a \enquote{literate execution} language
and leave model-written documentation aside, since in their system
\textcquote[sec.~7]{Bond2026LitExec}{@doc annotations are written by hand}.
\Textcite{Sotoudeh2025}, found through alphaXiv, propose literate tracing
and point out a limit of literate programming:
\textcquote[sec.~6.1]{Sotoudeh2025}{it cannot help make existing monolithic
codebases (GCC, Linux, etc.) more understandable}.

Four further records that the language model counted as connecting literate
programming with language models are not cited.
They are self-archived reports in the Zenodo repository: three were read in
abstract and contain no evaluation, and the fourth has no abstract and could
not be read.
They are listed in \cref{tab:lp-llm-direct}.

\section{Discussion and limitations}
\label{sec:ai-rq1-discussion}

The works that report gains do not settle whether literate prose helps.
The gains of verifiable literate programming come from a checking loop with
a simulated user who sees the answer, not from the prose alone, and those of
interoperable literate programming come from a small demonstration by the
tool's own developers.
The one ablation of literate comments found no difference in accuracy.
None of them measures what literate programming is for: whether people
understand the program.

Two things make \enquote{literate programming} mean something else in this
literature than it did for Knuth.
Where an evaluated system writes prose, a model generates it, mostly from the
code, after the fact,
whereas Knuth's programmer writes the explanation and orders the program for
it.
And \textcite{SunStaron2026} use \enquote{literate} for a property of the
model, which is partly a collision of terms.

The search has these limitations.
Semantic Scholar and DBLP could not be searched, and Scopus was not used.
The OpenAlex cells reached the limit of 100, and the coverage beyond it rests
on the uncapped title-and-abstract search, the citation chaining and
alphaXiv, which added grey literature but no peer-reviewed work in scope that
the queries had missed.
Knuth's own system, WEB, cannot be searched for as a keyword.
Two of the eleven works in scope were read only in abstract.
The language model's label for a supporting record means \enquote{connects
literate programming with language models explicitly}, not \enquote{shows a
benefit}, so the hit list's supporting rows are not evidence of benefit.
Six of the eleven works are preprints, and all but three are from 2025 or
2026, so the picture can change quickly.

\section{Conclusion}
\label{sec:ai-rq1-conclusion}

Since 2023, literate programming has been connected with language models
mainly through proposals and tools; only four works measure anything, and we
found none that tests literate programming in Knuth's sense, with
human-written prose and the program in expository order, against a
non-literate control.
The four measurements are narrow: a rating of model-written outlines by six
developers, gains in a simulation where the user knows the answer, a null
result for literate comments, and a study of model representations that uses
the word in another sense.
The lecture's answer to \cref{ex:ai-research-amount}, \enquote{mostly
proposals and tools, with little evaluation}, is the one this search
supports, and its statement that nobody has tested literate programming as
Knuth meant it holds for the databases and the citation trails searched
here.

\section{Further work}
\label{sec:ai-rq1-further}

This search left the following undone.
Semantic Scholar and DBLP should be searched when they answer, and Scopus
added.
The full texts of the two education papers\autocite{Birkenkrahe2024,
Rosiene2024} should be read before either is quoted beyond its abstract,
since one of them is the only counter-position in scope.

The field lacks, first, a controlled comparison of literate programs in
Knuth's sense with ordinary programs, when a model or an agent writes,
repairs or maintains the code.
It also lacks measurements of whether agents keep prose and code in step
over time, which the essay asserts\autocite{Whitlock2026} and the
environment paper doubts, having seen agents neglect such
artefacts\autocite{Burke2026notlob}; of what literate programs cost an
agent in tokens and time; of whether people understand literate
model-written programs better than ordinary ones; and of students using
literate programming with an AI assistant.
Every result so far comes from a single group, without replication.
If a controlled study found that literate programs help agents or their
readers, the lecture could state the benefit instead of the absence of
evidence; if it found no effect, the lecture's case for literate programming
would rest on the human reader alone.

\section{Hits bearing on the claim}
\label{sec:ai-rq1-hits}

\Cref{tab:lp-llm-direct} lists the records of the database searches that
bear on the question.
Knuth's article, the essay and the work on literate tracing came from the
known-document retrieval and the complementary searches and are not in it,
and the literate execution paper is classified as adjacent.

% Author's note (verification and reproduction):
%   scholar session : lp-llm-direct
%   report          : scratchpad lp-llm-direct-report.md (agent lit-rq1)
%   exports         : scratchpad lp-llm-direct.csv, lp-llm-direct.bib,
%                     lp-llm-direct-table.tex (copied here as
%                     litteratursokning/rq1-full.tex)
%   scratchpad      : /tmp/claude-1000/-home-dbosk-devel-edu-literate-
%                     programming/05aac7d0-554d-4ce5-a79b-906fad3fe64f/
%                     scratchpad/ (temporary: copy the exports into
%                     litteratursokning/ before it is cleaned)
%   searched        : 2026-11-05, 19:39-20:05 CET; providers openalex wos
%                     ieee arxiv hal, each separately, -y 2023- -l 100
%                     -f bibtex+prov, WoS as "TS=(<Q>)"
%                     (s2: key rejected, HTTP 403 -- failed every query;
%                     dblp: anti-bot page, HTTP 202 -- failed every
%                     query; scopus ok but not in the plan, not used)
%   superseded      : a first combined multi-provider run of S1 (one
%                     merged record, shared cap); session file moved to
%                     lp-llm-direct-work/superseded-session-combined.json;
%                     the redo recovered all its relevant results
%   complements     : OpenAlex works?filter=cites:W2342082360,
%                     from_publication_date:2023-01-01 (180 works);
%                     filter=title_and_abstract.search:"literate
%                     programming",from_publication_date:2023-01-01 (92);
%                     backward chaining of 8 papers; alphaXiv
%                     discover_papers x2; WebSearch x2; Knuth1984 read on
%                     the user's reMarkable (/Literate programming/Literate
%                     Programming, journal pp. 97 and 107)
%   classification  : scholar llm context (RQ1, four audit tags) + llm
%                     classify --no-examples --no-enrich, model
%                     openai-codex/gpt-5.6-sol, nothing pending; overrides
%                     via sessions decide: RossiPrestwich2026 and
%                     Rosiene2024 supports->qualifies, Yuan2026VLP
%                     qualifies->supports, Sun's 2026 thesis
%                     off-topic->adjacent; 11 rows < 0.7 read, kept
%                     adjacent
%   hit-list table  : scholar sessions export lp-llm-direct -f table
%                     --bearing-only --lang en (label tab:lp-llm-direct)
%   provenance      : ltnotes.bib, keys Knuth1984 SunStaron2026 Shi2025NLO
%                     Zhang2024ILP Yuan2026VLP RossiPrestwich2026
%                     Burke2026notlob VanClief2026ICM Birkenkrahe2024
%                     Rosiene2024 Whitlock2026 Bond2026LitExec Sotoudeh2025
%   dropped         : Starling & Dubinetska x2, LEAP, NALM (Zenodo grey);
%                     Sun 2026 thesis (not independent); Ustynov 2026 and
%                     Eghbali et al. 2026 (no LP link; used in RQ2)
% --- scholar LaTeX fragment --------------------------------------------
% Generated by: scholar sessions export -f table
% Session: lp-llm-direct
% \input this file from the body of a document whose
% preamble loads:
%   \usepackage{longtable}
%   \usepackage{booktabs}
%   \usepackage{array}
% ----------------------------------------------------------------------

\begingroup\footnotesize
\begin{longtable}{@{}%
  >{\raggedright\arraybackslash}p{0.44\textwidth}c%
  >{\raggedright\arraybackslash}p{0.13\textwidth}%
  >{\raggedright\arraybackslash}p{0.26\textwidth}@{}}
\caption{Records bearing on the claim (9 cited, 3 supporting and 1 qualifying or contradicting it, out of 567 unique hits; the 196 adjacent and 358 false hits are not listed). The last column gives the reason for each inclusion; for machine-classified rows the model's confidence is given. Databases: OpenAlex (OA; not open access), IEEE Xplore (IEEE), Web of Science (WoS), arXiv, HAL.}\label{tab:lp-llm-direct}\\
\toprule
Title & Year & Provider & Reason \\
\midrule
\endfirsthead
\multicolumn{4}{@{}l}{\emph{\tablename~\thetable{} (cont.)}}\\
\toprule
Title & Year & Provider & Reason \\
\midrule
\endhead
\midrule \multicolumn{4}{r@{}}{\emph{continued on next page}}\\
\endfoot
\bottomrule
\endlastfoot
A Literate Programming Environment for Human and Machine Agents & 2026 & arXiv & \emph{cited}: proposal or tool without evaluation \\
Grammar-Aware Literate Generative Mathematical Programming with Compiler-in-the-Loop & 2026 & OA & \emph{cited}: benchmark or computational evaluation \\
Guiding Human Validation of LLM-Generated Code via Verifiable Literate Programming & 2026 & arXiv & \emph{cited}: evaluation with human participants \\
Interpretable Context Methodology: Folder Structure as Agentic Architecture & 2026 & arXiv & \emph{cited}: proposal or tool without evaluation \\
Literate Programming With LLMs? - A Study on Rosetta Code and CodeNet & 2026 & WoS & \emph{cited}: benchmark or computational evaluation \\
Natural Language Outlines for Code: Literate Programming in the LLM Era & 2025 & WoS & \emph{cited}: evaluation with human participants \\
Renaissance of Literate Programming in the Era of LLMs: Enhancing LLM-Based Code Generation in Large-Scale Projects & 2024 & arXiv & \emph{cited}: benchmark or computational evaluation \\
THE ROLE OF AI CODING ASSISTANTS: REVISITING THE NEED FOR LITERATE PROGRAMMING IN COMPUTER AND DATA SCIENCE EDUCATION & 2024 & OA & \emph{cited}: position or practice paper \\
WIP: Focusing on Programmer Literacy in the Time of AI-Aided Code Generation & 2024 & WoS & \emph{cited}: position or practice paper \\
LEAP:Automating KPI Documentationin Enterprise Codebases through Literate Programming and AI & 2026 & OA & supports the claim (0.93) \\
Literate Programming: Donald Knuth, WEB, and Contemporary Workflows & 2026 & OA & supports the claim (0.99) \\
Prompt-Literate Workflow & 2026 & OA & supports the claim (0.99) \\
Towards a Non-Anthropocentric Logic Matrix (NALM): A Critical Review and Research Framing for AI-Native Intermediate Representations & 2026 & OA & qualifies or contradicts the claim (0.98) \\
\end{longtable}
\endgroup



% ============================================================= RQ2 + RQ3
\chapter{Does literate prose help language models, and does it help
  people check AI-written code?}
\label{app:ai-rq23}
\chapterprecis{The author has not yet reviewed the results in this
  appendix in full.}

The lecture makes two claims about prose next to code.
For the agent (\cref{sec:ai-agent-reader}), it claims that prose helps a
language model when it tells the model something correct that it couldn't
otherwise know, that wrong prose hurts more than none, and that prose
restating the code does little.
For the person who checks what an agent wrote (\cref{sec:ai-checker}), it
claims that explanations next to the code helped in some studies but not in
the one that varied how good they were, that comments make people adopt
AI-written code more often, and that AI assistance can lower understanding.

Literate prose here means any natural-language explanation kept with the
code: comments, docstrings, documentation, outlines and explanations.
Literate programs in Knuth's sense are a special case, and as
\cref{app:ai-rq1} found, nobody has tested them with language models, so
this chapter asks about the wider category.
The two research questions are:
\begin{description}
  \item[RQ2] Does prose of this kind help a language model or coding agent
    to write, repair or understand code, and when does it hurt?
  \item[RQ3] Does it help people to understand and validate code that an
    AI has written?
\end{description}
For RQ2, the claim is about conditions (when prose helps), so the search
must find both gains and losses.
For RQ3, the claim is that the evidence is mixed, which needs studies on
both sides.

\section{Method}
\label{sec:ai-rq23-method}

% Author's note (verification and reproduction): see the block above the
% hit-list placeholder at the end of the chapter.

The search was made on 5 November 2026 with \texttt{scholar}, in the same
databases as \cref{app:ai-rq1}: OpenAlex, Web of Science, IEEE Xplore, arXiv
and HAL worked; Semantic Scholar (access key rejected, HTTP 403) and DBLP
(a page that stops automated clients, HTTP 202) failed and were not
searched; Scopus was not part of the plan.
Each query was sent to each working database separately, with at most 50
hits per database, and to Web of Science as a topic search,
TS=(\dots), with the same keywords.

The support queries for RQ2 (S) and RQ3 (H) are:
\begin{description}
  \item[S1] (\enquote{code comments} OR docstrings OR documentation) AND
    (\enquote{large language model} OR \enquote{large language models} OR
    LLM) AND \enquote{code generation}
  \item[S2] \enquote{natural language outlines} AND (code OR LLM OR
    \enquote{large language model})
  \item[S3] (\enquote{code comments} OR docstrings OR documentation) AND
    (LLM OR \enquote{large language model} OR \enquote{large language
    models}) AND (\enquote{program repair} OR \enquote{code comprehension} OR
    \enquote{program comprehension} OR \enquote{code understanding})
  \item[S4] (documentation OR \enquote{code comments} OR docstrings OR
    \enquote{explanations of intent}) AND (\enquote{coding agent} OR
    \enquote{coding agents} OR \enquote{agentic coding} OR \enquote{software
    engineering agent})
  \item[S5] (\enquote{code readability} OR \enquote{readability} OR
    \enquote{self-documenting}) AND (LLM OR \enquote{large language model}
    OR \enquote{large language models}) AND (\enquote{code generation} OR
    \enquote{AI-generated code} OR \enquote{generated code})
  \item[S6] \enquote{comments} AND \enquote{code generation} AND
    (\enquote{large language model} OR LLM) AND (improve OR effect OR
    impact)
  \item[S7] (intent OR explanation OR \enquote{reasoning comments}) AND
    \enquote{code comments} AND (LLM OR \enquote{large language model})
  \item[H1] (\enquote{AI-generated code} OR \enquote{LLM-generated code} OR
    \enquote{AI-assisted programming}) AND (\enquote{program comprehension}
    OR \enquote{code comprehension} OR \enquote{code understanding}) AND
    (\enquote{controlled experiment} OR \enquote{human study} OR \enquote{user
    study})
  \item[H2] (\enquote{AI-generated code} OR \enquote{LLM-generated code})
    AND \enquote{code review} AND (explanation OR documentation OR comments
    OR \enquote{human study} OR \enquote{controlled experiment})
  \item[H3] (explanation OR explanations OR documentation OR
    \enquote{natural language}) AND (\enquote{AI-generated code} OR
    \enquote{LLM-generated code} OR \enquote{AI code}) AND (developers OR
    programmers) AND (understand OR comprehension OR validation OR
    verification)
  \item[H4] (\enquote{literate programming}) AND (LLM OR \enquote{large
    language model} OR \enquote{AI-generated} OR \enquote{LLM-generated})
    AND (validation OR understanding OR human)
  \item[H5] \enquote{human validation} AND (\enquote{LLM-generated code} OR
    \enquote{AI-generated code} OR \enquote{generated code})
  \item[H6] (\enquote{code summaries} OR \enquote{code summarization} OR
    \enquote{generated comments} OR \enquote{code explanations}) AND (LLM OR
    \enquote{AI-generated} OR \enquote{large language model}) AND
    (\enquote{controlled experiment} OR \enquote{user study} OR \enquote{human
    study}) AND (comprehension OR understanding)
\end{description}
The counter-queries for RQ2 (C) and RQ3 (D) look for misleading, outdated or
useless prose, its costs, and over-reliance on explanations:
\begin{description}
  \item[C1] (\enquote{misleading comments} OR \enquote{outdated comments} OR
    \enquote{comment inconsistency} OR \enquote{code-comment inconsistency}
    OR \enquote{incorrect comments}) AND (LLM OR \enquote{large language
    model} OR \enquote{large language models} OR \enquote{code generation})
  \item[C2] (\enquote{code comments} OR docstrings OR documentation) AND
    (LLM OR \enquote{large language model} OR \enquote{large language
    models}) AND (robustness OR perturbation OR adversarial OR misleading OR
    degrade OR \enquote{does not improve})
  \item[C3] (documentation OR \enquote{code comments} OR docstrings OR
    \enquote{context files}) AND (\enquote{coding agent} OR \enquote{coding
    agents} OR LLM) AND (\enquote{token cost} OR tokens OR overhead OR
    \enquote{no improvement})
  \item[C4] (\enquote{code readability} OR \enquote{self-documenting} OR
    \enquote{comment quality}) AND (LLM OR \enquote{large language model})
    AND (harm OR \enquote{no effect} OR limitation OR misleading OR
    \enquote{does not})
  \item[C5] (\enquote{comment removal} OR \enquote{without comments} OR
    \enquote{stripping comments} OR \enquote{removing comments} OR
    \enquote{comment-free}) AND (LLM OR \enquote{large language model} OR
    \enquote{large language models})
  \item[D1] (\enquote{over-reliance} OR overreliance OR \enquote{cognitive
    offloading} OR \enquote{comprehension debt} OR \enquote{cognitive debt})
    AND (\enquote{AI-generated code} OR \enquote{AI coding} OR \enquote{AI
    assistant} OR \enquote{code generation} OR LLM) AND (developers OR
    programmers OR \enquote{software engineering})
  \item[D2] (explanation OR explanations OR documentation) AND
    (\enquote{AI-generated code} OR \enquote{LLM-generated code} OR
    \enquote{AI code}) AND (overtrust OR \enquote{over-trust} OR
    \enquote{over-reliance} OR overreliance OR misleading OR wrong OR
    \enquote{does not help} OR incorrect)
  \item[D3] (\enquote{AI-generated code} OR \enquote{LLM-generated code} OR
    \enquote{AI-assisted}) AND (\enquote{code review} OR \enquote{reviewing})
    AND (burden OR effort OR \enquote{review time} OR fatigue OR \enquote{lost
    understanding} OR overreliance)
  \item[D4] (\enquote{illusion of understanding} OR \enquote{illusion of
    explanatory depth} OR \enquote{automation bias} OR \enquote{explanation})
    AND (\enquote{code generation} OR \enquote{AI-generated code} OR
    \enquote{programming assistant}) AND (\enquote{controlled experiment} OR
    \enquote{randomized} OR \enquote{human study} OR \enquote{user study})
  \item[D5] (\enquote{incorrect explanations} OR \enquote{wrong
    explanations} OR \enquote{misleading explanations} OR \enquote{inaccurate
    explanations} OR \enquote{inaccurate summaries} OR \enquote{misleading
    summaries}) AND (code OR programming) AND (LLM OR AI OR \enquote{large
    language model})
\end{description}
\Cref{tab:ai-rq23-queries} gives the hits.

\begin{table}[htbp]
  \begin{sidecaption}{Queries and number of hits per database, 5 November
    2026, with the tool \texttt{scholar}; the queries are given in
    \cref{sec:ai-rq23-method}.
    S = support query and C = counter-query on prose for language models
    (RQ2); H = support query and D = counter-query on prose for people
    (RQ3).
    Databases: OpenAlex (OA; not open access), Web of Science (WoS), IEEE
    Xplore (IEEE), arXiv, HAL (Hyper Articles en Ligne), Semantic Scholar
    (S2), DBLP.
    At most 50 hits were fetched per database and query, so 50 means
    \enquote{at least 50}.
    \enquote{--} means that the search failed: Semantic Scholar rejected the
    access key and DBLP answered every query with a page that stops
    automated clients, so neither database was searched.
    Web of Science got the keywords as \texttt{TS=(\dots)}, the others as
    boolean free text.}[tab:ai-rq23-queries]
    \centering\footnotesize
    \begin{tabular}{@{}lrrrrrrr@{}}
      \toprule
         & OA & WoS & IEEE & arXiv & HAL & S2 & DBLP \\
      \midrule
      S1 & 50 & 50  & 50   & 50    & 4   & -- & --   \\
      S2 & 49 & 1   & 0    & 1     & 0   & -- & --   \\
      S3 & 50 & 42  & 35   & 48    & 1   & -- & --   \\
      S4 & 50 & 10  & 6    & 50    & 0   & -- & --   \\
      S5 & 50 & 39  & 50   & 50    & 0   & -- & --   \\
      S6 & 50 & 10  & 18   & 41    & 0   & -- & --   \\
      S7 & 50 & 8   & 10   & 7     & 0   & -- & --   \\
      \midrule
      C1 & 49 & 8   & 7    & 9     & 0   & -- & --   \\
      C2 & 50 & 50  & 50   & 50    & 0   & -- & --   \\
      C3 & 50 & 44  & 47   & 50    & 0   & -- & --   \\
      C4 & 50 & 5   & 2    & 26    & 0   & -- & --   \\
      C5 & 50 & 4   & 2    & 8     & 1   & -- & --   \\
      \midrule
      H1 & 50 & 2   & 2    & 2     & 0   & -- & --   \\
      H2 & 50 & 1   & 4    & 7     & 0   & -- & --   \\
      H3 & 50 & 15  & 21   & 50    & 0   & -- & --   \\
      H4 & 50 & 2   & 1    & 4     & 0   & -- & --   \\
      H5 & 50 & 1   & 1    & 5     & 0   & -- & --   \\
      H6 & 50 & 3   & 3    & 3     & 0   & -- & --   \\
      \midrule
      D1 & 50 & 32  & 14   & 50    & 2   & -- & --   \\
      D2 & 50 & 2   & 1    & 9     & 0   & -- & --   \\
      D3 & 50 & 17  & 15   & 50    & 0   & -- & --   \\
      D4 & 50 & 1   & 2    & 4     & 0   & -- & --   \\
      D5 & 50 & 2   & 5    & 10    & 1   & -- & --   \\
      \bottomrule
    \end{tabular}
  \end{sidecaption}
\end{table}

HAL holds few records on this topic, so its zeros are answers, not failures.
The database searches were complemented the same day by two searches in
alphaXiv, for papers from 2023 onwards:
\begin{description}
  \item[Prose for the model] Keywords \enquote{code comments},
    \enquote{LLM}, \enquote{code generation}, \enquote{misleading
    comments}, \enquote{documentation} and \enquote{coding agents}, and the
    question \enquote{Do natural-language code comments or documentation
    improve or harm LLM performance in code generation, repair,
    comprehension, or agentic tasks?}
  \item[Prose for people] Keywords \enquote{AI-generated code},
    \enquote{program comprehension}, \enquote{code review},
    \enquote{over-reliance}, \enquote{cognitive offloading},
    \enquote{comprehension debt} and \enquote{explanation}, and the question
    \enquote{Controlled experiments with human developers: do explanations,
    comments or documentation help people understand, review or validate
    AI-generated code, or do they increase over-reliance and lost
    understanding?}
\end{description}
Two of the included papers were found only through alphaXiv, and three were
found by both alphaXiv and the queries.
Two papers were pointed out by the parallel search of \cref{app:ai-rq1}.
One of them\autocite{eghbali2026equivalence} had already been returned by
the queries S2, S3 and H4, and was read after the pointer.
The other\autocite{ustynov2026beyond}, on the formats of log files for
agents, was returned by no query and was retrieved by its arXiv identifier
as a known document.
Before the search, five papers known to be relevant were set aside to test
the queries, and the queries found all five.

The 2,436 hits came to 1,666 unique records.
Their titles were filtered by keywords for each question, about 180 titles
were read, about 60 papers were opened in full text or abstract, and 34 were
included; with the two papers pointed out by the other search, 36 sources
are cited, 19 for RQ2 and 18 for RQ3, one of them for both.
Several were read in the abstract and introduction only, and the
chapter says so where it matters.
A language model is classifying every record of the search in the same way
as for the other chapters, but that classification was not finished when
this chapter was written, so the screening reported here is the keyword
screening alone.

\section{Results}
\label{sec:ai-rq23-results}

\subsection{Prose for the language model}

\paragraph{What the prose says matters, not that there is prose.}
The cleanest causal study, by \textcite{pan2026talking} and peer-reviewed,
controlled what comments said, and found that
\textcquote[abstract]{pan2026talking}{Comments from source solutions that
pass the tests raise recipient pass@1 by 17.2\,\% on average.
In contrast, comments describing failed solutions provide no reliable gain,
while comments written for a different problem reduce pass@1 by
20.8\,\%.}
The 17.2 is in percentage points, the recipients were four small open
models, and the comments came from a stronger model's passing solution, so
they partly pass on the solution itself.
Asking a model to comment its own code helped much less:
\textcquote[sec.~1]{pan2026talking}{the best prompt recovers at most about
24\,\% of the external-comment lift}.

\paragraph{With the code present, more documentation did not help.}
For an agent fixing issues, \textcquote[abstract]{arman2026compact}{When the
source is present, neither static compact documentation nor retrieved
context beats the issue alone}; the full-length description did worst, and
only when the source was withheld did documentation raise the share of
tests passed, from 0.08 to 0.71.
Repository context files for agents likewise gave no reliable gain, at more
than 20\,\% higher cost\autocite{gloaguen2026agentsmd}; \cref{app:ai-rq4}
treats them in detail.
Documentation can also replace code the model does not see: docstrings that
describe a function exactly let a model predict its output with
\textcquote[abstract]{eghbali2026equivalence}{12.8-24.5\,\% higher
accuracy, when compared to human-written documentation and documentation
generated by a baseline approach}, in a setting where the model was shown the
docstring and not the code.

\paragraph{Comments help where they carry requirements.}
In automated bug fixing, \textcquote[abstract]{vitale2026comments}{comments
improve ABF accuracy by up to threefold when present in both phases}, though
from low levels, for example from 2.78 to 9.80\,\% exact fixes.
In translation between programming languages, model-written comments
changed success by between $-90$ and $+435$\,\%\autocite{gupta2026translation},
a nearly balanced mix of gains and losses.
On one legacy module, comments raised four models' quiz scores from 84 to
90--96\,\%, and wholly wrong comments lowered them to
61\,\%\autocite{sabetto2025legacy}; the models scored 75\,\% without the code
at all, so the module may be in their training data.
Comments that users could edit raised the share of solved tasks of a code
model by 17.1\,\%, and users of the tool were 16.7\,\% faster, in a study with
12 participants\autocite{di2025mutual}.
Switching the concept of a comment on and off inside the model shifted
performance between $-90$ and $+67$\,\%, depending on model and
task\autocite{imani2026insideout}.

\paragraph{Wrong prose hurts.}
\Textcite{macke2024documentation} found that
\textcquote[abstract]{macke2024documentation}{providing an LLM with
\enquote{incorrect} documentation can greatly hinder code understanding, while
incomplete or missing documentation does not seem to significantly affect an
LLM's ability to understand code}.
Misleading cues in natural language lowered 17 models' output prediction by
23.2\,\% on average\autocite{lam2025codecrash}.
Defective code left in comments made coding assistants write defective code
in up to 58.17\,\% of cases\autocite{huang2026commenttraps}, and insecure
instructions hidden in comments gave vulnerable completions in 77.4--92.3\,\%
of cases, against 1.7--5.1\,\% without them\autocite{munir2026hidden}.
One study points the other way, for detecting vulnerabilities rather than
writing code: adversarial comments had no significant effect on eight models,
and removing the comments hurt the weaker models\autocite{%
thornton2026adversarial}.

\paragraph{Prose costs tokens, and removing it is not free either.}
Docstrings in prompts could be shortened by 25--40\,\% without lowering the
quality of the generated code\autocite{yang2026lessismore}, whereas
compressing the log files an agent reads lowered its input by 17\,\% but
raised the cost of the session by 67\,\%, in one model and one session per
format\autocite{ustynov2026beyond}.
In real agent sessions, agent instruction files and working notes made up
60.5\,\% of the agents' interactions with documentation, and consulting them
was associated with less testing\autocite{gao2026agentfriendly}.

\paragraph{Literate prose itself.}
The only work in this search that uses literate programming by name is the
outline study already discussed in \cref{app:ai-rq1}, which shows that a
model can write such prose, not that it helps\autocite{Shi2025NLO}.
A survey of practitioners by a tool vendor found that 81\,\% still consider
readability important for AI-written code\autocite{%
takerngsaksiri2025readability}, which is an opinion, not a test.

\subsection{Prose for the person who checks}

\paragraph{Tool studies with correct explanations are positive.}
With explanations placed next to generated code, participants, mostly
graduate students,
\textcquote[sec.~5.4]{yan2024ivie}{answered an average of 90.2\,\% of
questions correctly, in contrast to 65.0\,\% with the baseline}, and the
baseline was a chatbot.
Linked natural-language summaries raised task correctness from 81.7 to
95.0\,\% for 20 developers, though the authors warn that their results
\textcquote[limitations]{tang2026naturaledit}{are also subject to
automation bias: a readable NL diff may increase perceived control while
still concealing semantic errors or inaccurate NL-code
mappings}.
Explanations from a language model helped 32 participants understand code
more than web search did, for code written by people\autocite{nam2024gilt}.
The study of editable comments\autocite{di2025mutual} and the verifiable
literate programming study\autocite{Yuan2026VLP} point the same way, with 12
and eight participants.
None of these studies varied whether the explanation was correct.

\paragraph{Varying the quality of the explanation removed the gain.}
In a pre-registered experiment with 86 developers judging assertions about
AI-generated functions, \textcite{kaufman2026assertions} found that
\textcquote[sec.~RQ2]{kaufman2026assertions}{the presence of a comment (even
if that comment was of high-quality) did not significantly improve
participant accuracy compared to providing no comment at all (p > 0.1)}.
Under-specified comments lowered accuracy compared with exact ones (odds
ratio 0.58) while raising confidence compared with no comment.
The stimuli were assertions about small benchmark functions, not whole
programs.

\paragraph{Explanations raise adoption and trust, not necessarily
correctness.}
In an experiment with 173 programmers, \textcquote[abstract]{%
li2026comments}{the presence of comments significantly influences
programmers' adoption of AI-generated code regardless of the participants'
development expertise}, which measures uptake, not whether the code was
right.
Full explanations of AI code-review recommendations gave the highest
trust, but not the highest agreement, and the recommendations were always
correct, since \textcquote[threats]{gao2026xai}{Evaluating whether
explanations help developers detect incorrect recommendations would require
a different design}.
In a study of students working with an agent, comprehension fell
sharply, and more comments in the code were associated with lower
comprehension, a correlation rather than a manipulated
effect\autocite{balepur2026impaired}.
Outside programming, \textcquote[abstract]{kim2025reliance}{the presence of
explanations increases reliance on both correct and incorrect responses},
which has not been tested for code.

\paragraph{AI assistance can lower understanding.}
In a randomized study by Anthropic researchers,
\textcquote[sec.~1.1]{shen2026skill}{using AI assistance to complete tasks
that involve this new library resulted in a reduction in the evaluation
score by 17\,\% or two grade points (Cohen's d = 0.738, p = 0.010)}.
With an AI assistant, 15 participants performed better
but \textcquote[abstract]{qiao2025faster}{showed no corresponding
improvement in overall comprehension (p = 0.59)}.
Beginners understanding code from a language model succeeded in 32.5\,\% of
tasks\autocite{zi2025beginners}, and labelling code as written by a
language model changed where reviewers looked, not how thoroughly they
reviewed\autocite{khojah2026scrutiny}.

\paragraph{Debt and reviews.}
A review of practitioners' reports finds that
\textcquote[results]{tavantzis2026cognitive}{a plausible AI-generated
explanation does not necessarily produce durable human understanding}, and
an essay names \enquote{intent debt}, a
\textcquote[introduction]{storey2026debt}{lack of externalized goals,
constraints, and rationale that both humans and AI systems need to work
safely and efficiently with the codebase}.
Students' diaries show AI code accepted as a black box and verification
bypassed\autocite{ahmad2026comprehension}, and a systematic review of
generative AI in programming education reports that its assistants
\textcquote[abstract]{qiao2026slr}{often yield inaccurate or unclear
explanations}.

\section{Discussion and limitations}
\label{sec:ai-rq23-discussion}

For the model, the split between studies that find gains and studies that
find none follows what the prose adds.
Comments help when they carry requirements or intent the model cannot infer,
as in bug fixing, translation and legacy code, and do nothing when they
restate code already in context or a well-known benchmark function; the
causal study shows this directly, since the content and not the commenting
carried the effect\autocite{pan2026talking}.
Documentation that replaces code the model cannot see helps a
lot\autocite{eghbali2026equivalence,arman2026compact}, but literate prose
is an addition to the code, not a replacement, so the regime with the code
present is the one that applies to literate programming.
Wrong prose is followed when a model writes code, while detection with
strong models seems more robust; the counter-example is a single preprint.

For people, the positive studies compared explanations with a chatbot, web
search or other tools, kept the explanations correct, and measured
understanding by quizzes or self-reports right after the task.
The null result varied the quality of the explanation and measured whether
people judged correctly.
These designs answer different questions, so the studies do not contradict
each other; but only the second tests the situation a checker is in, where
the explanation might be wrong.
That comments increase adoption fits both readings: people use commented
code more, not necessarily more correctly.
Some authors have interests in the result: the outline study and one of the
explanation studies have Google authors, the readability survey evaluates its
authors' own tool, and the study of lost understanding is by Anthropic
researchers, against the interest of an AI company.

The search has these limitations.
Semantic Scholar and DBLP could not be searched, and Scopus was not used.
Each cell was limited to 50 hits, which OpenAlex reached for all but two
queries, so relevant older items may be buried by the relevance ranking.
The screening of 1,666 titles was assisted by keywords rather than a full
reading, and the language model's classification of the whole search was not
finished, so the hit list is not yet included.
Several cited papers were read in abstract and introduction only.
Nearly all of the work is benchmarks on models from 2024--2026 and can date
quickly, and none of it uses a literate program in Knuth's sense.

\section{Conclusion}
\label{sec:ai-rq23-conclusion}

Prose helps a language model when it adds correct information the model
doesn't otherwise have, hurts it when it is wrong, and does little but cost
tokens when it restates code the model can see; for people checking
AI-written code, the evidence is mixed, and the one experiment that varied
the quality of the explanation found no gain and a cost in false
confidence.
Our confidence is moderate for the first answer, which several independent
benchmarks support, and low to moderate for the second, which rests on few
studies with people.
Neither answer is about literate programming in Knuth's sense, which nobody
has tested with language models or with people checking AI-written code.
The lecture's claims in \cref{sec:ai-agent-reader,sec:ai-checker} are stated
at this strength, with the reservation that they are about prose with code in
general.

\section{Further work}
\label{sec:ai-rq23-further}

This search left the following undone.
The language model's classification of the whole search must be finished and
its hit list added to this chapter.
Semantic Scholar and DBLP should be searched when they answer, and Scopus
added.
Several papers were found but could not be obtained in full text, the most
important first: a study of outdated comments in LLM-based code completion
(Xia, Huang, Chen and Zheng, ASE 2026), a neurophysiological study of code
understanding in modern coding workflows (ISSRE 2025), an evaluation of the
explanations language models give for the patches they write (Sobania and
others, SSBSE 2023), a conceptual framework for comprehension gaps in
AI-assisted development (IEEE Access 2026), and a study of how programmers'
AI literacy affects their handling of bad code (BISE 2026).
% DOIs of the papers to fetch: 10.1145/3832783.3834359 (When Comments
% Mislead, ACM 403), 10.1109/ISSRE66568.2025.00058,
% 10.1007/978-3-031-48796-5_12, 10.1109/ACCESS.2026.3689574,
% 10.1007/s12599-026-01000-5; also 10.1145/3597503.3639176 (TODO comments
% and Copilot, ICSE 2024; abstract read, not cited).

The field lacks, first, a controlled study of literate programs, with the
chunks in explanatory order, either with language models or with people who
check AI-written code.
It also lacks an experiment with people that crosses whether an explanation
is present with whether it is correct, for whole AI-written programs; one
study did so for assertions\autocite{kaufman2026assertions}, and another left
that design to future work\autocite{gao2026xai}.
Nobody has compared prose written by a person with prose written by a model,
for the same code; the causal benchmark suggests that a model's own comments
add little for the model\autocite{pan2026talking}, and nothing tests the same
for people.
Nor has anyone followed over time whether prose reduces the debt in
understanding that practitioners report, or how stale prose affects agents.
If such studies found that human-written literate prose helps checkers find
errors, the lecture could recommend literate programs to anyone who checks an
agent's work; if they found the same null result as for comments, the
lecture's caution about explanations would extend to literate programs too.

\section{Hits bearing on the claim}
\label{sec:ai-rq23-hits}

The classification of this search's records was still running when this
appendix was written, so its hit list is not yet included.

% Author's note (verification and reproduction):
%   scholar session : lp-prose-effects
%   report          : scratchpad lp-prose-effects-report.md (lit-rq23)
%   exports         : scratchpad lp-prose-effects/lp-prose-effects.bib
%                     (36 entries); per-query results in
%                     lp-prose-effects/q/, matrix.log, merged.json
%   searched        : 2026-11-05, ~19:40-20:10 CET; providers openalex
%                     ieee arxiv hal (-l 50 -f csv) and wos as
%                     TS=(<same keywords>), each separately, -n
%                     lp-prose-effects (s2: key rejected, HTTP 403 --
%                     failed every query; dblp: bot wall, HTTP 202 --
%                     failed every query; scopus ok, not in the plan)
%   superseded      : a first test query ("code comments" AND "large
%                     language model" AND "code generation", five
%                     providers in one call, no TS= for WoS, 57 hits),
%                     replaced by S1; its records are among the session's
%                     1,800 but not in the matrix
%   complements     : alphaXiv discover_papers x2 (retained 2601.20245
%                     shen2026skill and 2604.13277 ahmad2026comprehension
%                     only there); alphaXiv answer_pdf_queries on "When
%                     Comments Mislead" (not indexed); Crossref/OpenAlex
%                     metadata; arXiv ti: queries for four closed papers
%                     (no copies); hand-offs from team lead: Eghbali
%                     (already in S2, S3, H4) and Ustynov 2604.07502
%                     (known item by id)
%   classification  : scholar llm context (RQ2/RQ3) + llm classify
%                     lp-prose-effects --no-examples --no-enrich, model
%                     openai-codex/gpt-5.6-sol, batches of 100; STILL
%                     RUNNING at 20:12 (620 of 1,800 done, see
%                     lp-prose-effects/classify.log); theme tags on the
%                     cited papers not yet applied
%   hit-list table  : TODO, when classification is done:
%                     scholar sessions export lp-prose-effects -f table
%                     --bearing-only --lang en --label tab:lp-prose-effects
%                     --theme ... -o litteratursokning/rq23-full
%                     then uncomment the \input below
%   provenance      : ltnotes.bib, keys pan2026talking arman2026compact
%                     vitale2026comments macke2024documentation
%                     sabetto2025legacy gupta2026translation
%                     imani2026insideout lam2025codecrash
%                     huang2026commenttraps munir2026hidden
%                     thornton2026adversarial yang2026lessismore
%                     di2025mutual gao2026agentfriendly
%                     takerngsaksiri2025readability kaufman2026assertions
%                     balepur2026impaired shen2026skill gao2026xai
%                     yan2024ivie nam2024gilt tang2026naturaledit
%                     li2026comments qiao2025faster zi2025beginners
%                     khojah2026scrutiny kim2025reliance
%                     tavantzis2026cognitive storey2026debt
%                     ahmad2026comprehension qiao2026slr
%                     eghbali2026equivalence ustynov2026beyond
%                     gloaguen2026agentsmd Shi2025NLO Yuan2026VLP
% TODO: table pending classification
% \input{litteratursokning/rq23-full}


% ================================================================ RQ4
\chapter{Do context files, skills and spec-driven development help coding
  agents?}
\label{app:ai-rq4}
\chapterprecis{The author has not yet reviewed the results in this
  appendix in full.}

The lecture makes several claims about prose that instructs an agent rather
than explains the code (\cref{sec:ai-instructions}).
Context files, such as \texttt{AGENTS.md} or \texttt{CLAUDE.md}, did not
significantly change whether agents solved their tasks, and made them more
expensive.
Out-of-date files are followed over the code, even when the agent notices the
conflict.
Skills written by people helped on specialized tasks and those the agent
wrote itself made things worse, while on ordinary repositories most public
skills had no effect.
When rating these approaches (\cref{ex:ai-rate-instructions}), the lecture
adds that spec-driven development, where a specification is written before
the code, has hardly been evaluated beyond its authors' own projects, and when
rating Showboat (\cref{ex:ai-rate-showboat}) that no scientific evaluation of
such documents was found.

The research question is: do context files, skills and spec-driven
development help coding agents, and what do the studies measure as help?
Help can mean that the agent solves more tasks, that it costs less, or that
its work is easier for a person to check, and the chapter keeps these apart.
The claims tested are mostly negative ones (no significant gain, no
evaluation), which need studies that could have found a gain, and searches
for evidence in both directions.

\section{Method}
\label{sec:ai-rq4-method}

% Author's note (verification and reproduction): see the block above the
% hit-list \input at the end of the chapter.

The search was made on 5 November 2026 with \texttt{scholar}, in the same
databases as \cref{app:ai-rq1,app:ai-rq23}: OpenAlex, Web of Science, IEEE
Xplore, arXiv and HAL worked; Semantic Scholar (access key rejected, HTTP
403) and DBLP (a page that stops automated clients, HTTP 202) failed and were
not searched; Scopus was not part of the plan.
Each query was sent with at most 40 hits per database.
Each item gives the query as sent to OpenAlex, IEEE Xplore, arXiv and HAL,
then the same keywords as sent to Web of Science.
The support queries are:
\begin{description}
  \item[S1] \enquote{AGENTS.md} OR \enquote{CLAUDE.md} coding agent; WoS:
    TS=((\enquote{AGENTS.md} OR \enquote{CLAUDE.md}) AND coding AND agent*)
  \item[S2] \enquote{context file} OR \enquote{context files} coding agent
    evaluation; WoS: TS=((\enquote{context file} OR \enquote{context files})
    AND coding AND agent* AND evaluat*)
  \item[S3] \enquote{agent instructions} OR \enquote{repository
    instructions} OR \enquote{copilot instructions} OR \enquote{cursor
    rules} empirical; WoS: TS=((\enquote{agent instructions} OR
    \enquote{repository instructions} OR \enquote{copilot instructions} OR
    \enquote{cursor rules}) AND empirical)
  \item[S4] \enquote{agentic coding} instructions effect efficiency impact;
    WoS: TS=(\enquote{agentic coding} AND instruction* AND (effect* OR
    efficien* OR impact*))
  \item[S5] \enquote{agent skills} coding agent evaluation; WoS:
    TS=(\enquote{agent skills} AND coding AND agent* AND evaluat*)
  \item[S6] \enquote{spec-driven development} OR
    \enquote{specification-driven development} agent; WoS:
    TS=((\enquote{spec-driven development} OR
    \enquote{specification-driven development}) AND agent*)
  \item[S7] executable demonstration proof of agent work coding agent
    verify; WoS: TS=((executable NEAR/3 (demonstration* OR documentation))
    AND coding AND agent*)
  \item[S8] \enquote{instruction files} OR \enquote{agent manifests} OR
    \enquote{agent configuration files} coding agent; WoS:
    TS=((\enquote{instruction files} OR \enquote{agent manifests} OR
    \enquote{agent configuration files}) AND coding AND agent*)
  \item[S9] architecture documentation prompt context coding agent
    experiment developers experience capability; WoS: TS=(architecture AND
    documentation AND context AND \enquote{coding agent*} AND experiment*)
  \item[S10] agent proof of work demo verify reproducible artifact coding
    agent evidence of correctness supervisor; WoS: TS=((\enquote{proof of
    work} OR \enquote{executable demo*} OR \enquote{reproducible evidence})
    AND \enquote{coding agent*})
  \item[S11] \enquote{rule files} OR \enquote{persistent configuration} OR
    \enquote{AI configuration} coding agents; WoS: TS=((\enquote{rule files}
    OR \enquote{persistent configuration} OR \enquote{AI configuration}) AND
    coding AND agent*)
  \item[S12] \enquote{skill files} OR \enquote{agent skills} coding agent
    cost success controlled; WoS: TS=((\enquote{skill files} OR
    \enquote{agent skills}) AND agent* AND (cost OR success))
  \item[S13] instruction following coding agents instruction surfaces rules
    project files; WoS: TS=(\enquote{instruction following} AND
    \enquote{coding agent*})
  \item[S14] coding agent adoption code quality technical debt
    configuration repositories; WoS: TS=(\enquote{coding agent*} AND
    adoption AND (\enquote{code quality} OR \enquote{technical debt}))
\end{description}
The counter-queries reuse the concept terms with words for harm, cost,
staleness, conflict and mistakes:
\begin{description}
  \item[C1] \enquote{AGENTS.md} OR \enquote{CLAUDE.md} OR \enquote{context
    files} coding agent reduced success cost tokens harmful; WoS:
    TS=((\enquote{AGENTS.md} OR \enquote{CLAUDE.md} OR \enquote{context
    files}) AND coding AND agent* AND (reduc* OR cost OR token* OR harm*))
  \item[C2] \enquote{agent instructions} OR \enquote{repository
    instructions} OR \enquote{cursor rules} stale outdated redundant; WoS:
    TS=((\enquote{agent instructions} OR \enquote{repository instructions}
    OR \enquote{cursor rules}) AND (stale OR outdated OR redundant))
  \item[C3] \enquote{spec-driven development} OR
    \enquote{specification-driven development} limitations does not help
    cost empirical; WoS: TS=((\enquote{spec-driven development} OR
    \enquote{specification-driven development}) AND (limitation* OR cost OR
    empirical))
  \item[C4] \enquote{agent skills} coding agent no benefit worse security
    limitations; WoS: TS=(\enquote{agent skills} AND coding AND agent* AND
    (benefit* OR worse OR limitation* OR security))
  \item[C5] \enquote{instruction files} OR \enquote{context files} stale
    outdated agent worse performance not helpful; WoS:
    TS=((\enquote{instruction files} OR \enquote{context files}) AND (stale
    OR outdated OR worse) AND agent*)
  \item[C6] coding agent documentation code disagree conflict instruction
    files coherence; WoS: TS=(\enquote{coding agent*} AND (documentation OR
    \enquote{instruction files}) AND (conflict* OR disagree* OR coheren*))
  \item[C7] \enquote{AGENTS.md} OR \enquote{CLAUDE.md} smells mistakes
    problems configuration; WoS: TS=((\enquote{AGENTS.md} OR
    \enquote{CLAUDE.md}) AND (smell* OR mistake* OR problem*))
\end{description}
\Cref{tab:ai-rq4-queries} gives the hits.

\begin{table}[htbp]
  \begin{sidecaption}{Queries and number of hits per database, 5 November
    2026, with the tool \texttt{scholar}; the queries are given in
    \cref{sec:ai-rq4-method}.
    S = support query, C = counter-query.
    Databases: OpenAlex (OA; not open access), IEEE Xplore (IEEE), arXiv,
    HAL (Hyper Articles en Ligne), Web of Science (WoS), Semantic Scholar
    (S2), DBLP.
    \enquote{Union} is the number of unique records from OpenAlex, IEEE
    Xplore, arXiv and HAL together, as the search session stored them.
    At most 40 hits were fetched per database and query, so 40 means
    \enquote{at least 40}.
    \enquote{--} means that the search failed: Semantic Scholar rejected the
    access key and DBLP answered every query with a page that stops
    automated clients, so neither database was searched.
    Web of Science got the keywords as \texttt{TS=(\dots)}, the others as
    free text.}[tab:ai-rq4-queries]
    \centering\footnotesize
    \begin{tabular}{@{}lrrrrrrrr@{}}
      \toprule
          & OA & IEEE & arXiv & HAL & Union & WoS & S2 & DBLP \\
      \midrule
      S1  & 40 & 4    & 40    & 40  & 121   & 2   & -- & --   \\
      S2  & 40 & 32   & 40    & 14  & 122   & 0   & -- & --   \\
      S3  & 40 & 2    & 40    & 39  & 119   & 0   & -- & --   \\
      S4  & 40 & 0    & 40    & 0   & 74    & 0   & -- & --   \\
      S5  & 40 & 0    & 40    & 0   & 68    & 1   & -- & --   \\
      S6  & 40 & 3    & 40    & 40  & 120   & 2   & -- & --   \\
      S7  & 40 & 1    & 40    & 0   & 80    & 1   & -- & --   \\
      S8  & 40 & 7    & 40    & 40  & 125   & 3   & -- & --   \\
      S9  & 40 & 0    & 40    & 0   & 80    & 0   & -- & --   \\
      S10 & 40 & 0    & 40    & 0   & 80    & 0   & -- & --   \\
      S11 & 40 & 11   & 40    & 40  & 130   & 0   & -- & --   \\
      S12 & 40 & 1    & 40    & 0   & 70    & 7   & -- & --   \\
      S13 & 40 & 0    & 40    & 0   & 75    & 2   & -- & --   \\
      S14 & 40 & 0    & 40    & 0   & 77    & 1   & -- & --   \\
      \midrule
      C1  & 40 & 4    & 40    & 0   & 83    & 1   & -- & --   \\
      C2  & 40 & 1    & 40    & 0   & 80    & 0   & -- & --   \\
      C3  & 40 & 2    & 40    & 4   & 85    & 2   & -- & --   \\
      C4  & 40 & 0    & 40    & 0   & 80    & 0   & -- & --   \\
      C5  & 39 & 6    & 40    & 0   & 85    & 0   & -- & --   \\
      C6  & 40 & 0    & 40    & 0   & 79    & 0   & -- & --   \\
      C7  & 40 & 4    & 40    & 0   & 82    & 0   & -- & --   \\
      \bottomrule
    \end{tabular}
  \end{sidecaption}
\end{table}

The search session stored one merged record set per query, so the counts per
database were taken by running the same queries again outside the session.
Web of Science indexes almost none of this literature, which is mostly
preprints from 2026, so its zeros carry little weight, and the HAL cells of
40 are mostly false hits on single words of the query.

Seven papers known to be relevant were set aside before the search, and the
queries S1--S10 and C1--C5 found all seven.
The reference list of one of the central studies\autocite{xu2026stale}
showed six relevant papers that the queries had missed.
The queries S11--S14, C6 and C7 add the terms the field uses for these
things, \enquote{rule files}, \enquote{persistent configuration},
\enquote{AI configuration}, \enquote{instruction following},
\enquote{smells} and conflict or coherence, and with them all six were
found.
Three further papers, all descriptive, were still missed by every query; they
are named in \cref{sec:ai-rq4-further}.

Searches outside the databases, all on the same day, served mostly to read
the sources and establish their status:
\begin{description}
  \item[alphaXiv] A search with the keywords \enquote{coding agents},
    \enquote{proof of work}, \enquote{executable documentation} and
    \enquote{verification evidence} and the question \enquote{Research on
    coding agents producing executable demonstrations or verifiable evidence
    (proof) of their work for a human supervisor, and whether agents
    misreport or fabricate what they did} returned 13 papers, of which two
    were kept\autocite{smyth2026overclaiming,xu2026ctxgap}.
    The full texts of four papers were read through alphaXiv.
  \item[Practitioner sources] Simon Willison's blog post on Showboat and
    Rodney\autocite{willison2026showboat} and the Showboat
    documentation\autocite{willison2026showboatreadme} were retrieved as
    known documents and read in full, together with the open issue about
    agents that edit the document directly.
  \item[The web] One web search, \enquote{\enquote{Evaluating AGENTS.md}
    Gloaguen ICLR 2026 workshop openreview}, led to the authors' page that
    gives the venue of one paper\autocite{gloaguen2026agentsmd}.
\end{description}

The roughly 2,300 hits came to 1,377 unique records.
A language model classified all of them, as supporting, qualifying or
contradicting the claim, adjacent, or false hits.
Every record it called supporting or qualifying, and the one record it was
unsure of, was read again: 41 of its calls were changed to adjacent, mostly
papers that propose a way to optimize skills rather than test whether skills
help, and those changes rest on title and abstract.
About 205 records were read in title or abstract and about 60 in abstract or
more.
Thirty-one sources are included: 21 read in full text, partly (the abstract
and the passages that bear on the question), eight in abstract only, and two
practitioner sources.

\section{Results}
\label{sec:ai-rq4-results}

\subsection{Context files: controlled experiments}

The most careful benchmark so far ran four agents on 438 tasks, one sample per
task, and \textcite{gloaguen2026agentsmd} report that
\textcquote[abstract]{gloaguen2026agentsmd}{providing context files does not
generally improve task success rates, while increasing inference cost by
over 20\,\% on average}.
Files written by the developers raised success by 2.4 percentage points,
which is not significant, and files written by a model lowered it slightly;
the overview of the repository, the part such files most often contain, did
not help the agents find the relevant files:
\textcquote[sec.~4.3]{gloaguen2026agentsmd}{We conclude that context files
are not effective at providing a repository overview}.
The paper is a workshop paper, and its first version said that context files
tend to reduce success; the later version quoted here says they make no
significant difference.

Three smaller experiments agree.
With two agents on 17 tasks, \textcquote[abstract]{khatri2026context}{%
Context strategy does not measurably move correctness on either agent}, though
the author notes that the study can detect only large effects.
On 58 tasks chosen to be neither too easy nor too hard, rule files raised
success by 7 to 14 percentage points, but \textcquote[results]{%
zhang2026guardrails}{no condition is significantly different from any other},
and \textcquote[sec.~1]{zhang2026guardrails}{Random rules improve a coding
agent's task performance as much as expert-curated ones}.
In 1,650 sessions, the size, position and structure of the file and
contradictions in it had no detectable effect on whether the agent followed
it\autocite{mcmillan2026adherence}, read in abstract only.

Two experiments vary what the file contains.
Including every category of instruction \textcquote[abstract]{%
polandra2026categories}{does not reliably reduce agent-file-generation cost
or improve generated code quality}, in a task far from everyday work: having
an agent regenerate a whole library.
The format of an architecture specification hardly mattered for the strongest
models, \textcquote[abstract]{canedo2026architecture}{On weaker models, format
produces spreads of 0.83--2.42 points}, in a single-author preprint scored
partly by a language model.

One experiment finds lower cost: with \texttt{AGENTS.md},
\textcquote[abstract]{lulla2026agentsmd}{the presence of AGENTS.md is
associated with a lower median runtime (\(\Delta\)28.64\,\%) and reduced output
token consumption (\(\Delta\)16.58\,\%)}, for one agent on small pull
requests.
Total tokens did not change significantly, and correctness was not measured,
since \textcquote[sec.~3.1.8]{lulla2026agentsmd}{A comprehensive evaluation
of the output quality \textelp{} is beyond the scope of this paper}.

\subsection{Context files in real projects}

Two observational studies of real repositories point in no clear direction.
After instruction files were added to 148 projects,
\textcquote[abstract]{arabat2026instructions}{27.7\,\% of the projects
increased their merge rate by at least 20\,\%, while 26.35\,\% decreased
it}.
In 441 repositories, those that adopted agents without committed AI
configuration saw complexity rise faster, and the authors \textcquote[%
abstract]{denisovblanch2026markdown}{present these findings as
hypothesis-generating}.
Two descriptive studies find that the files \textcquote[abstract]{%
chatlatanagulchai2026readmes}{evolve like configuration code through
frequent, small additions}, and that problems such as rules better left to a
linter and bloated context are common\autocite{dossantos2026smells}.

\subsection{Out-of-date files and instruction following}

The clearest result is a harmful one.
In a controlled experiment with 40 projects, \textcquote[abstract]{%
xu2026stale}{Without the sentence, 90.1\,\% of the agents' choices followed
the stale file.
Claude Code pointed out the conflict in 79 of 80 runs and still followed the
file.
Codex CLI rewrote the existing code to match the file in 20 of 80 runs, and
the project's own tests then failed in 17 of them.}
The sentence is one instruction to prefer the code when the two disagree.
A second experiment found the same: \textcquote[abstract]{%
mohammadi2026workingset}{where standard and code disagree, agents follow the
standard even when it prescribes the worse code, so a stale convention file
costs more than no file}.
In a case study of one repository, stale documentation was the largest class
of failures, 48 of 118\autocite{aliabadi2026lessons}, and an existing
consistency checker found stale references to code in 23.0\,\% of 356
repositories\autocite{treude2026contextrot}.

\Textcite{chakrabarti2026growing} explain why the files keep growing:
\textcquote[abstract]{chakrabarti2026growing}{once an instruction's rationale
is gone, deleting it without risking a correctness regression costs
\(O(2^{|D|})\) in a prompt of \(|D|\) instructions}.
Comments with the rationale removed 99.3\,\% of the excess instructions, but
that was tested on instruction-following tasks, not on coding agents.

Two benchmarks qualify what following a rule means.
\textcquote[abstract]{huang2026harnessif}{When a coding agent obeys a rule, it
may simply have been going to do that anyway}; every model was worse at rules
that go against its prior, by 3.6 to 7.4 points.
And patches that pass the tests still \textcquote[abstract]{truong2026swecc}{%
violate 43.1 percent of applicable project policies}, read in abstract only.

\subsection{Skills}

The results on skills split by the kind of task.
On specialized tasks built together with their skills,
\textcquote[abstract]{li2026skillsbench}{Curated Skills raise the average
pass rate from 33.9\,\% to 50.5\,\% (+16.6 percentage points; 25.5\,\%
normalized gain)}, but \textcquote[results]{li2026skillsbench}{Self-generated
Skills land below the no-Skills baseline on all three configurations}, by 8.1
to 11.5 percentage points.
On real software repositories, \textcquote[abstract]{han2026sweskills}{92 of
98 skills have no statistically significant effect, and adding skills raises
the average pass ratio only from 76.6\,\% to 76.9\,\%}, and harmful skills
override the task's own instructions.

Four further studies qualify the gains.
Skills helped in all nine configurations of one benchmark, but
\textcquote[sec.~1]{cheng2026utility}{The same Skills yield gains under some
configurations and losses under others on 36.78\,\% of tasks}.
A study of failures found that \textcquote[abstract]{dong2026harmful}{%
seemingly relevant skills often make the agent incorrectly implement or omit
task-required implementation elements}.
Shortening skills by up to 54\,\% did not lower the cost, since the skill is a
small and mostly cached part of what the agent reads\autocite{%
xu2026shortening}.
And skills go stale as software changes: \textcquote[abstract]{%
duan2026repo2skill}{every evaluated transition invalidates part of the V1
skill set}, read in abstract only.

\subsection{Spec-driven development}

The first systematic study of spec-driven development as practised combines
interviews,
online discussions and a corpus of specification files, and notes that
\textcquote[sec.~1]{agarwal2026sdd}{Tool documentation, gray literature, and
early research prescribe workflows and promise better quality, less agent
drift, and higher productivity.
However, their evaluations rarely extend beyond the authors' own projects}.
The practitioners \textcquote[abstract]{agarwal2026sdd}{used SDD selectively
and mainly for new features, and reported more reliable agent behavior
alongside greater planning and review costs}, and they reported that the
specifications decayed; the benefits are
self-reported.
The only measurement of outcomes is of tokens: cutting a full specification
to a bare user story \textcquote[abstract]{smekal2026taskspec}{raises token
spend by 29.7\,\%}, for one model, without measuring correctness, read in
abstract only.

\subsection{Proof of agent work}

Showboat is described only in practitioner sources.
Its author wants \textcquote{willison2026showboat}{tools that allow agents to
clearly demonstrate their work to me, while minimizing the opportunities for
them to cheat about what they've done}, and its documentation says that
\textcquote{willison2026showboatreadme}{A verifier can re-execute all code
blocks and confirm the outputs still match}.
We found no scientific evaluation of documents like these.
The closest is a pilot of 64 runs, in which models rather than people did the
reviewing and \textcquote[abstract]{schmalbach2026contracts}{delegation
contracts bought reviewability rather than correctness}.
The problem such documents address has been measured, in abstracts only:
\textcquote[abstract]{smyth2026overclaiming}{agents fail to read every file
they were asked to review in 67.9\,\% of runs}, and readers of an agent's
reports alone recover 52.0\,\% of the task state, against 92.7\,\% with the
full context\autocite{xu2026ctxgap}.

\section{Discussion and limitations}
\label{sec:ai-rq4-discussion}

Most of the apparent conflicts are different measurements.
The study that found lower cost measured time and output tokens for one agent
on small tasks, while the one that found higher cost measured money, which
input and reasoning tokens dominate, for four agents.
The two are not fully reconciled, since the second also found a higher cost
for the agent the first one used.
The two benchmarks of skills differ in their tasks: specialized procedures
with a low baseline in one, ordinary repositories and a high baseline in the
other.
Both agree that skills help when they supply procedures the agent would
otherwise leave out.
That random rules did as well as curated ones points to priming rather than
content.
The observational studies are confounded by the discipline of the teams and
by time, as their authors say.

Taken together, the studies suggest that text helps an agent when it carries
information the agent can't get from the code or its training: procedures in
curated skills, facts that have been withheld, structured specifications for
weak models.
Overviews that restate the code did not help.
That reading is ours, and no single study tests it.

Some authors have interests in the result, though none of the studies
evaluates its own company's product: authors at Amazon, at ByteDance and at a
company that builds coding agents appear among the experiments, and one
author of the spec-driven development study works at Google, which makes one
of the tools in that study's corpus.

The search has these limitations.
Semantic Scholar and DBLP could not be searched, and Scopus was not used.
The limit of 40 hits was reached by arXiv for every query and by OpenAlex
for all but one.
The counts per database were taken outside the session, and the session's
records do not say which database returned them, so the provider column of
the hit list names every database the query was sent to.
Several full-text readings covered the abstract and the passages that bear on
the question rather than the whole paper, eight sources were read in abstract
only, and the language model's classification was changed on title and
abstract.
Nearly everything is a 2026 preprint or workshop paper, on single model
versions and mostly on Python, and the experiments ran between 17 and 438
tasks each.

\section{Conclusion}
\label{sec:ai-rq4-conclusion}

Measured by whether the agent solves its task, repository context files have
not been shown to help, and out-of-date ones demonstrably harm, since agents
follow them over the code; skills written by people help on specialized tasks
but mostly not on ordinary repositories, and those agents write for
themselves hurt; spec-driven development and executable demonstrations like
Showboat's have no evaluation of outcomes.
Our confidence is moderate that generic context files give no gain in success
and that stale files harm, and low for their effect on cost, which depends on
the setup.
The lecture's claims in \cref{sec:ai-instructions} and its ratings in
\cref{ex:ai-rate-instructions,ex:ai-rate-showboat} are stated at this
strength.

\section{Further work}
\label{sec:ai-rq4-further}

This search left the following undone.
Semantic Scholar and DBLP should be searched when they answer, and Scopus
added.
Three descriptive papers were missed by every query and should be read: on
the availability and granularity of skills (Xu and Wu, 2026), on context
engineering for AI agents in open-source software (Mohsenimofidi and others,
MSR 2026), and on the context developers provide (Jiang and Nam).
The sources read only in abstract should be read in full before any of them
carries a claim, the experiment on file structure\autocite{%
mcmillan2026adherence} and the pilot on reviewability\autocite{%
schmalbach2026contracts} first.
% arXiv ids: 2605.31408 (Xu & Wu; downloaded, not read), 2512.18925 (Jiang
% & Nam).  Also to read: han2026sweskills results by configuration,
% gloaguen2026agentsmd appendix B, li2026skillsbench appendix D.6, the
% xu2026stale data (github.com/WenJing95/stale-docs-data).

The field lacks, first, a test of prose that explains \emph{why} and is
interleaved with the code, as in literate programming, against a separate
instruction file; the closest are comparisons of specification
formats\autocite{canedo2026architecture} and of overviews\autocite{%
gloaguen2026agentsmd}, and rationale comments tested outside
coding\autocite{chakrabarti2026growing}.
Nobody has compared spec-driven development with working without a
specification on real tasks, for success, defects or cost, and nobody has
evaluated executable demonstration documents or whether re-running them stops
fabricated output.
There is no study of developers with and without context files, no study of
whether keeping the files current pays off over time, and none of long work
over many sessions, where persistent context should matter most; the
experiments that exist are single tasks.
If an experiment found that prose interleaved with the code goes stale less
and helps agents more than a separate file, the lecture could recommend
literate programs as the place for an agent's instructions; if not, the
rating of instruction files in \cref{ex:ai-rate-instructions} would stand as
it is.

\section{Hits bearing on the claim}
\label{sec:ai-rq4-hits}

\Cref{tab:lp-agent-context} lists the records that bear on the question.
The practitioner sources on Showboat and the two papers found through
alphaXiv are not in it.

% Author's note (verification and reproduction):
%   scholar session : lp-agent-context (42 search records, 1,377 unique)
%   report          : scratchpad lp-agent-context-report.md (lit-rq4)
%   exports         : scratchpad lp-agent-context.csv, lp-agent-context.bib
%                     (31 entries), lp-agent-context-table.tex (copied
%                     here as litteratursokning/rq4-full.tex); per-provider
%                     split lp-agent-context-work/matrix2.tsv
%   searched        : 2026-11-05, 19:40-19:54 CET; S1-S10 and C1-C5 at
%                     19:40-19:44, S11-S14, C6, C7 at 19:53-19:54;
%                     providers openalex ieee arxiv hal --limit 40, wos as
%                     TS=(...) (s2: key rejected, HTTP 403 -- failed every
%                     query; dblp: bot wall, HTTP 202 -- failed every
%                     query; scopus ok, not in the plan)
%   provider column : the hit list's provider column shows the union of
%                     the providers queried, since the session stores one
%                     merged record set per query; the split per provider
%                     was recomputed outside the session
%   recall          : 7 brief items found by S1-S10/C1-C5; 6 gaps found by
%                     chaining xu2026stale (zhang2026guardrails,
%                     denisovblanch2026markdown, mohammadi2026workingset,
%                     dossantos2026smells, huang2026harnessif,
%                     xu2026shortening) closed by concept queries S11-S14,
%                     C6, C7; still missed: 2605.31408, Mohsenimofidi et al.
%                     MSR'26, 2512.18925
%   complements     : alphaXiv get_paper_content (2602.11988, 2601.20404,
%                     2607.27250, 2610.30571); alphaXiv discover_papers
%                     (retained 2609.20812, 2610.20003); WebFetch/curl of
%                     simonwillison.net/2026/Feb/10/showboat-and-rodney/,
%                     github.com/simonw/showboat, issue #12; WebSearch
%                     for the Gloaguen venue; Crossref and arXiv API
%   classification  : scholar llm classify --no-examples, model
%                     openai-codex/gpt-5.6-sol, 1,377 records, none
%                     pending; all supports/qualifies rows read, plus the
%                     one row < 0.7 (discarded); 41 overrides to
%                     adjacent-subtopic; theme tags ctx-effect ctx-stale
%                     ctx-following ctx-descriptive skills sdd
%                     proof-of-work on the cited sources
%   hit-list table  : scholar sessions export lp-agent-context -f table
%                     --bearing-only --lang en (label tab:lp-agent-context)
%   provenance      : ltnotes.bib, keys gloaguen2026agentsmd
%                     lulla2026agentsmd khatri2026context
%                     zhang2026guardrails mcmillan2026adherence
%                     polandra2026categories canedo2026architecture
%                     xu2026stale mohammadi2026workingset
%                     aliabadi2026lessons treude2026contextrot
%                     chakrabarti2026growing huang2026harnessif
%                     truong2026swecc arabat2026instructions
%                     denisovblanch2026markdown chatlatanagulchai2026readmes
%                     dossantos2026smells li2026skillsbench han2026sweskills
%                     cheng2026utility dong2026harmful xu2026shortening
%                     duan2026repo2skill agarwal2026sdd smekal2026taskspec
%                     schmalbach2026contracts smyth2026overclaiming
%                     xu2026ctxgap willison2026showboat
%                     willison2026showboatreadme
% --- scholar LaTeX fragment --------------------------------------------
% Generated by: scholar sessions export -f table
% Session: lp-agent-context
% \input this file from the body of a document whose
% preamble loads:
%   \usepackage{longtable}
%   \usepackage{booktabs}
%   \usepackage{array}
% ----------------------------------------------------------------------

\begingroup\footnotesize
\begin{longtable}{@{}%
  >{\raggedright\arraybackslash}p{0.44\textwidth}c%
  >{\raggedright\arraybackslash}p{0.13\textwidth}%
  >{\raggedright\arraybackslash}p{0.26\textwidth}@{}}
\caption{Records bearing on the claim, Agent context files, skills and SDD (27 cited, 14 supporting and 22 qualifying or contradicting it, out of 1361 unique hits; the 593 adjacent and 705 false hits are not listed). The last column gives the reason for each inclusion; for machine-classified rows the model's confidence is given. Databases: OpenAlex (OA; not open access), IEEE Xplore (IEEE), Web of Science (WoS), arXiv, HAL.}\label{tab:lp-agent-context}\\
\toprule
Title & Year & Provider & Reason \\
\midrule
\endfirsthead
\multicolumn{4}{@{}l}{\emph{\tablename~\thetable{} (cont.)}}\\
\toprule
Title & Year & Provider & Reason \\
\midrule
\endhead
\midrule \multicolumn{4}{r@{}}{\emph{continued on next page}}\\
\endfoot
\bottomrule
\endlastfoot
A Few Pages of Markdown: Committed AI Configuration and Lower Quality Cost after Coding-Agent Adoption & 2026 & OA, IEEE, arXiv, HAL & \emph{cited}: effect of context files \\
Agent READMEs: An Empirical Study of Context Files for Agentic Coding & 2026 & OA, IEEE, arXiv, HAL & \emph{cited}: descriptive studies \\
Agent Skills Can Be Harmful: An Empirical Study of Skill-Induced Failures in LLM Agents & 2026 & OA, IEEE, arXiv, HAL & \emph{cited}: agent skills \\
An Empirical Study of Agent Skills' Downstream Utility & 2026 & OA, IEEE, arXiv, HAL & \emph{cited}: agent skills \\
Architecture as Capability Equalizer for Coding Agents & 2026 & OA, IEEE, arXiv, HAL & \emph{cited}: effect of context files \\
Can your AI agent be cheaper? Investigating the effects of task specifications on token spend in agentic coding tasks & 2026 & OA, IEEE, arXiv, HAL & \emph{cited}: spec-driven development \\
Configuration Smells in AGENTS.md Files: Common Mistakes in Configuring Coding Agents & 2026 & OA, IEEE, arXiv, HAL & \emph{cited}: descriptive studies \\
Context Rot in AI-Assisted Software Development: Repurposing Documentation Consistency for AI Configuration Artifacts & 2026 & OA, IEEE, arXiv, HAL & \emph{cited}: staleness and maintenance \\
Correct Code, Broken Contributions? SWE-CC: Benchmarking Repository Policy Compliance for Coding Agents & 2026 & OA, IEEE, arXiv, HAL & \emph{cited}: instruction following \\
Do Context Files Help Coding Agents? A Two-Agent Ablation Study on Real Repositories & 2026 & OA, IEEE, arXiv, HAL & \emph{cited}: effect of context files \\
Evaluating AGENTS.md: Are Repository-Level Context Files Helpful for Coding Agents? & 2026 & OA, IEEE, arXiv, HAL & \emph{cited}: effect of context files \\
Guardrails Beat Guidance: A Large-Scale Study of Rules, Skills, and Persistent Configuration for Coding Agents & 2026 & OA, IEEE, arXiv, HAL & \emph{cited}: effect of context files \\
Harness-IF: Evaluating Instruction Following Across Instruction Surfaces in Coding Agents & 2026 & OA, IEEE, arXiv, HAL & \emph{cited}: instruction following \\
Instruction Adherence in Coding Agent Configuration Files: A Factorial Study of Four File-Structure Variables & 2026 & OA, IEEE, arXiv, HAL & \emph{cited}: instruction following \\
Lessons Outlive Their Failures: A Case Study of Rules for AI Coding Agents & 2026 & OA, IEEE, arXiv, HAL & \emph{cited}: staleness and maintenance \\
On the Impact of AGENTS.md Files on the Efficiency of AI Coding Agents & 2026 & OA, IEEE, arXiv, HAL & \emph{cited}: effect of context files \\
Repo2Skill-Evo: Repository Skills Go Stale in Silence & 2026 & OA, IEEE, arXiv, HAL & \emph{cited}: staleness and maintenance \\
Shortening agent skills leaves LLM agent cost unchanged & 2026 & OA, IEEE, arXiv, HAL & \emph{cited}: agent skills \\
SkillsBench: Benchmarking How Well Agent Skills Work Across Diverse Tasks & 2026 & OA, IEEE, arXiv, HAL & \emph{cited}: agent skills \\
Software Delegation Contracts: Measuring Reviewability in AI Coding-Agent Work & 2026 & OA, IEEE, arXiv, HAL & \emph{cited}: proof of agent work \\
Spec-Driven Development as Prescribed and as Practiced & 2026 & OA, IEEE, arXiv, HAL & \emph{cited}: spec-driven development \\
Stale Instruction Files Override the Code in Coding Agents: A One-Sentence Fix & 2026 & OA, IEEE, arXiv, HAL & \emph{cited}: staleness and maintenance \\
SWE-Skills-Bench: Do Agent Skills Actually Help in Real-World Software Engineering? & 2026 & OA, IEEE, arXiv, HAL & \emph{cited}: agent skills \\
The Role of Agentic Instruction Categories for Code Generation: A Study of Five Python Libraries & 2026 & OA, IEEE, arXiv, HAL & \emph{cited}: effect of context files \\
The Working Set of a Coding Agent: Coherence Debt in Repository-Scale Tasks & 2026 & OA, IEEE, arXiv, HAL & \emph{cited}: staleness and maintenance \\
Toward Instructions-as-Code: Understanding the Impact of Instruction Files on Agentic Pull Requests & 2026 & OA, IEEE, arXiv, HAL & \emph{cited}: effect of context files \\
Why Does CLAUDE.md Keep Growing? Catastrophic Remembering in Agentic Coding & 2026 & OA, IEEE, arXiv, HAL & \emph{cited}: staleness and maintenance \\
A Lean and Spec-Driven AI-Assisted Software Development Lifecycle for Applied AI Education: The AI-SDLC Approach & 2026 & OA, IEEE, arXiv, HAL & supports the claim (0.98) \\
Agent Skill Evolution: How Revisions Affect Coding Agents & 2026 & OA, IEEE, arXiv, HAL & supports the claim (1.00) \\
Bootstrapping Coding Agents: The Specification Is the Program & 2026 & OA, IEEE, arXiv, HAL & supports the claim (0.98) \\
Codified Context: Infrastructure for AI Agents in a Complex Codebase & 2026 & OA, IEEE, arXiv, HAL & supports the claim (0.94) \\
Constitutional Spec-Driven Development: Enforcing Security by Construction in AI-Assisted Code Generation & 2026 & OA, IEEE, arXiv, HAL & supports the claim (0.97) \\
Context Engineering for Multi-Agent LLM Code Assistants Using Elicit, NotebookLM, ChatGPT, and Claude Code & 2025 & OA, IEEE, arXiv, HAL & supports the claim (0.94) \\
ContextCov: Deriving and Enforcing Executable Constraints from Agent Instruction Files & 2026 & OA, IEEE, arXiv, HAL & supports the claim (0.96) \\
Enhancing Specification-Driven Development with LLM-Generated Transition Diagrams: An Empirical Evaluation & 2026 & OA, IEEE, arXiv, HAL & supports the claim (1.00) \\
Hardening vibe-coded web applications with an automated security-audit skill: a controlled comparison of two builds of the same app & 2026 & OA, IEEE, arXiv, HAL & supports the claim (0.99) \\
Small Language Models and Spec-Driven Development for High-Accuracy Agentic AI Systems & 2026 & OA, IEEE, arXiv, HAL & supports the claim (0.98) \\
SPECIFICATION-DRIVEN DEVELOPMENT AS STRUCTURED APPROACH FOR SOURCE CODE GENERATION & 2026 & OA, IEEE, arXiv, HAL & supports the claim (0.98) \\
Specification-Driven Governance for Fairness-by-Design in Agentic Software Engineering: How Far Can We Go? & 2026 & OA, IEEE, arXiv, HAL & supports the claim (0.97) \\
Use of GitHub Copilot Agent Mode for Automated Development in R\&D Data Engineering & 2026 & OA, IEEE, arXiv, HAL & supports the claim (0.98) \\
Zoro: Active Controls to Support Rule Following in Vibe Coding & 2026 & OA, IEEE, arXiv, HAL & supports the claim (1.00) \\
A Framework for Evaluating Agentic Skills at Scale & 2026 & OA, IEEE, arXiv, HAL & qualifies or contradicts the claim (0.99) \\
Can Coding Agents Migrate to Post-Quantum Cryptography? & 2025 & OA, IEEE, arXiv, HAL & qualifies or contradicts the claim (0.94) \\
Citation Discipline in Spec-Driven Development: A Cross-Model Empirical Study of Output Determinism and Automated Hallucination Detection in LLM-Generated Code & 2026 & OA, IEEE, arXiv, HAL & qualifies or contradicts the claim (0.99) \\
CLAUDE.md and AGENTS.md size: what instruction files cost in Claude Code and Codex & 2026 & OA, IEEE, arXiv, HAL & qualifies or contradicts the claim (0.99) \\
Continuous Process-Level Evaluation for Evolving Enterprise AI Agent Skills & 2026 & OA, IEEE, arXiv, HAL & qualifies or contradicts the claim (0.99) \\
Curating Always-Loaded Context for LLM Agents: A Capacitated Assortment Model with Censored Feedback & 2026 & OA, IEEE, arXiv, HAL & qualifies or contradicts the claim (0.99) \\
Developer-on-the-Loop: Engineering Cognitive Systems with Coding Agents & 2026 & OA, IEEE, arXiv, HAL & qualifies or contradicts the claim (0.90) \\
Evaluating Agent Skills for Version-Specific Plugin Migration: A Retrospective Study & 2026 & OA, IEEE, arXiv, HAL & qualifies or contradicts the claim (0.99) \\
M4Bench: Evaluating Procedural Specification for Clinical EHR Derivation Agents & 2026 & OA, IEEE, arXiv, HAL & qualifies or contradicts the claim (1.00) \\
OctoBench: Benchmarking Scaffold-Aware Instruction Following in Repository-Grounded Agentic Coding & 2026 & OA, IEEE, arXiv, HAL, WoS & qualifies or contradicts the claim (0.98) \\
One Skill Takes Over: Clashing Agent Skills and a One-File Fix & 2026 & OA, IEEE, arXiv, HAL & qualifies or contradicts the claim (1.00) \\
Position, Role or Label? A Controlled Study of Instruction Conflicts in LLM Agents & 2026 & OA, IEEE, arXiv, HAL & qualifies or contradicts the claim (0.99) \\
Practical Implementation Report on Introducing Spec-Driven Development Using AI Agents in Software Development PBL & 2026 & OA, IEEE, arXiv, HAL & qualifies or contradicts the claim (0.98) \\
Report: Progressive Disclosure of Agent Skills & 2026 & OA, IEEE, arXiv, HAL & qualifies or contradicts the claim (0.99) \\
Security Impact of Coding Agent's Memory Files & 2026 & OA, IEEE, arXiv, HAL & qualifies or contradicts the claim (1.00) \\
Skill Constellations: Tracing the Supply Chain of Agent Skills on GitHub & 2026 & OA, IEEE, arXiv, HAL & qualifies or contradicts the claim (0.99) \\
Subagents vs Agent Skills: Executing Reusable Knowledge for Long-Horizon Agentic Tasks & 2026 & OA, IEEE, arXiv, HAL & qualifies or contradicts the claim (1.00) \\
What Is a Skill Worth? Structure-Aware Shapley Valuation of Agent Skills & 2026 & OA, IEEE, arXiv, HAL & qualifies or contradicts the claim (0.94) \\
What Should a Skill Remember? Quality--Cost Trade-offs in Cost-Aware Skill Rewriting for Language Model Agents & 2026 & OA, IEEE, arXiv, HAL & qualifies or contradicts the claim (1.00) \\
When \textquotedbl{}Do Not\textquotedbl{} Is Not Deny: Security Rules in CLAUDE.md vs Built-In Controls & 2026 & OA, IEEE, arXiv, HAL & qualifies or contradicts the claim (1.00) \\
When and How Context Rot Appears in Coding Agents: A White-Box Study of Agent Skills in Code Auditing & 2026 & OA, IEEE, arXiv, HAL & qualifies or contradicts the claim (0.99) \\
Zero Human Code: Building Operational Software with AI Coding Agents Without Code Inspection & 2026 & OA, IEEE, arXiv, HAL & qualifies or contradicts the claim (0.98) \\
\end{longtable}
\endgroup

